\documentclass[11pt]{article}
\usepackage[margin=25mm]{geometry}
\usepackage{fontspec}
\setmainfont{Latin Modern Roman}
\usepackage{amsmath,amssymb,mathtools}
\usepackage{unicode-math}
\setmathfont{Latin Modern Math}
\usepackage{xurl,hyperref}
\hypersetup{colorlinks=true,urlcolor=blue}
\setcounter{secnumdepth}{0}
\setlength{\parindent}{0pt}
\setlength{\parskip}{5pt}
\setlength{\emergencystretch}{3em}
\newcommand{\sref}[2]{\hyperlink{source:#1:#2}{[S#2]}}
\newcommand{\cataloguescope}[1]{\paragraph{Recorded target.}#1}
\newcommand{\R}{\mathbb R}
\newcommand{\C}{\mathbb C}
\newcommand{\N}{\mathbb N}
\newcommand{\Z}{\mathbb Z}
\newcommand{\Q}{\mathbb Q}
\newcommand{\F}{\mathbb F}
\newcommand{\eps}{\varepsilon}
\begin{document}
\section{Functional Differential Equations for Many-Body Green Functions}\label{statement}

\subsection{Definitions and mathematical statement}

\paragraph{Shared notation.}$\mathbb N=\{1,2,\ldots\}$; the standard number fields have their usual meanings. All additional conventions are stated locally.


\paragraph{Definitions and conventions.}
An interacting one-body Green function is a time-ordered or otherwise conventionally ordered two-point correlation function of creation and annihilation fields. The Dyson equation writes $G=G_0+G_0\Sigma G$, with convolution in the space, spin and time variables and self-energy $\Sigma$. A functional differential equation treats $G$ as a functional of an added external potential and differentiates with respect to that potential, thereby encoding higher correlations without assuming a known self-energy.

\paragraph{Target retained from the canonical source catalogue.}
\paragraph{Catalogue wording retained for traceability.}
How can the interacting Green functions of many-body quantum systems be obtained from a functional differential equation rather than the usual Dyson integral equation?

\paragraph{Source scope and qualification.}
The source\textquotesingle s ordering, contour and external-field convention fix the functional equation. Producing formal perturbative derivatives alone does not construct the desired nonperturbative interacting Green function.
\subsection{Short English statement}
Obtain interacting many-body Green functions directly from their equations in functional derivatives.
\subsection{Research attempt}
\paragraph{Model, date and investigated scope.}
OpenAI GPT-6 Astra (configured model \texttt{gpt-6-astra}), ultra reasoning; generated 11 October 2026. Derived and audited an exact source-differential equation for a quartic scalar integral; exposed missing physical initial data rather than constructing the many-body Green function.

\paragraph{Sources and literature check.}
Read the supplied definition, catalogue wording, source qualifications and reference entries for this question. The historical source/context and inherited status are retained; no fresh literature or novelty certification is claimed. No external theorem is used without proof in this attempt; explicit model assumptions are distinguished from facts about the source model.

\paragraph{A nonperturbative scalar analogue.} For real \(J\), \(a>0\), \(\lambda>0\), define
\[
 Z(J)=\int_{\mathbb R}\exp(-ax^2/2-\lambda x^4/4+Jx)\,dx.
\]
The quartic decay dominates every polynomial times the integrand on compact \(J\)-intervals, justifying differentiation and integration by parts. Integrating the derivative in \(x\) gives the exact equation
\[
 \lambda Z^{(3)}(J)+aZ'(J)-JZ(J)=0.
\]
Since \(Z>0\), its normalized one-point function \(g=Z'/Z\) obeys
\[
 \lambda(g''+3gg'+g^3)+ag-J=0.
\]
This follows by twice differentiating \(Z'=gZ\). Thus an external-source differential equation can encode interaction nonperturbatively without inserting a self-energy.
\paragraph{The selection gap.} The linear third-order equation requires three initial values; evenness gives \(Z'(0)=0\), and normalization fixes one scale, but \(Z''(0)/Z(0)\) still depends on the physical integral. The equation alone does not determine it. As \(\lambda\to0\), the equation loses order and becomes \(aZ'=JZ\), so perturbative continuity is a singular selection problem rather than an arbitrary initial-value prescription.
\paragraph{Source scope.} This finite scalar integral is an elementary test of the method, not the source's ordered fermionic two-point functional equation. Its contour, operator ordering, physical boundary conditions and infinitely many source variables have not been reconstructed. The exact first gap is to prescribe sufficient physical functional boundary data with existence and uniqueness for the actual many-body equation. No Dyson-free construction of that Green function is claimed.

\paragraph{Verification and adversarial audit.}
Justified differentiation by quartic domination, vanishing integration boundaries, nonlinear substitution and the singular lambda=0 order reduction. This is a GPT-6 Astra Ultra self-audit, not independent human verification. Prompt A drafting followed by prompt B scope/dependency/boundary self-audit. Justified differentiation by quartic domination, vanishing integration boundaries, nonlinear substitution and the singular lambda=0 order reduction. The remaining gap is: Specify and solve the physical functional boundary-value problem with the source contour and ordering. Unresolved status and zero original-target completion are retained.

Coordinator checked the integration-by-parts sign, logarithmic-derivative equation and independent even initial datum. The scalar model is not the fermionic functional equation. This additional check was model-assisted, not independent human verification.

\paragraph{Tools, process and contributions.}
Python 3 and local file-reading/writing tools for manuscript serialization and structural checks; mathematical calculations were symbolic reasoning, with no numerical experiment or theorem prover. Contribution prompt A was followed by prompt B. GPT-6 Astra with ultra reasoning drafted the research and performed its stated self-audit; any additional coordinator contribution is identified in the verification and audit. A GPT-6-based Codex coordinator assembled the manuscript and checked repository integration; its exact variant and reasoning configuration were not exposed. No human mathematical contribution, independent human proof check or formal proof certification is supplied. Complete derivations and the consulted source references are included here; private reasoning transcripts are not published.

\paragraph{Final status and first remaining gap.}
\textbf{UNRESOLVED}. Specify and solve the physical functional boundary-value problem with the source contour and ordering.

\section{Completion Estimate}
\noindent\textbf{COMPLETION: 0\%}
The proved calculation is an elementary auxiliary result or a test of a restricted strategy; no new portion of the full source target is established.

\subsection{Sources}
\begin{itemize}
\item[\hypertarget{source:statement:1}{S1}] ULAM.AI and the UnsolvedMath Contributors, UnsolvedMath, catalogue entry. Catalogue curation under CC BY 4.0; definitions and English exposition added locally.
\url{https://huggingface.co/datasets/ulamai/UnsolvedMath}.
\item[\hypertarget{source:statement:2}{S2}] Mathematical Methods in Quantum Chemistry (2018), Oberwolfach Reports, DOI 10.4171/owr/2018/13, pp. 631--721. Primary report accessed 2026-10-11; relevant PDF pages 56.
\url{https://doi.org/10.4171/owr/2018/13}.
\item[\hypertarget{source:statement:3}{S3}] Bold Feynman diagrams (catalogue-listed context; not independently consulted in this preparation).
\url{https://arxiv.org/abs/1809.02901}.
\item[\hypertarget{source:statement:4}{S4}] Coupled-cluster Green functions (catalogue-listed context; not independently consulted in this preparation).
\url{https://pmc.ncbi.nlm.nih.gov/articles/PMC8493088/}.
\item[\hypertarget{source:statement:5}{S5}] Green-function implementation (catalogue-listed context; not independently consulted in this preparation).
\url{https://doi.org/10.1016/j.cpc.2024.109380}.
\end{itemize}
\label{end:statement}
\end{document}
